% Con il continuo avanzamento delle tecnologie di calcolo quantistico e con il crescente interesse verso i collegamenti di comunicazione quantistica per l'interconnessione di sistemi quantistici remoti, la distribuzione affidabile di stati quantistici attraverso fibre ottiche sta diventando un requisito tecnologico fondamentale.

% Le coppie di fotoni entangled in polarizzazione sono particolarmente adatte ai collegamenti in fibra, poiché la polarizzazione può essere manipolata e misurata mediante componenti ottici standard. Le fibre reali, tuttavia, introducono trasformazioni coerenti e meccanismi di decoerenza che modificano le correlazioni osservate e possono degradare l'entanglement distribuito. Questa tesi sviluppa un framework progressivo che collega tali effetti fisici a previsioni numeriche e a quantità accessibili sperimentalmente.

% La modellizzazione parte da una descrizione, motivata sperimentalmente, dell'evoluzione residua della polarizzazione dopo una compensazione parziale. Viene inizialmente introdotta una fase relativa fissa per distinguere le variazioni coerenti delle correlazioni misurate da un'effettiva degradazione dell'entanglement. Il modello viene quindi esteso a fluttuazioni di fase non risolte, riassunte dal parametro di coerenza \( \mu=\langle e^{i\theta}\rangle \), il cui modulo determina, all'interno di questo modello, la concurrence residua attraverso \( C=|\mu| \). Successivamente vengono introdotti effetti depolarizzanti indipendenti nei due rami in fibra tramite il fattore di sopravvivenza delle correlazioni \( \eta=(1-p_A)(1-p_B) \), e combinati con le fluttuazioni di fase in una descrizione comune mediante operatore densità. Infine, i modelli efficaci vengono collegati alla dispersione modale di polarizzazione, nella quale fotoni a banda finita subiscono trasformazioni birifrangenti dipendenti dalla frequenza e uno stato di polarizzazione misto emerge dopo aver effettuato la traccia sui gradi di libertà spettrali non risolti.

% Il framework è implementato in una libreria MATLAB modulare e standalone che riproduce numericamente questa gerarchia di modellizzazione. Il simulatore fornisce una pipeline comune per la preparazione dello stato, la composizione dei canali, la media statistica, l'analisi delle correlazioni di polarizzazione e la valutazione di purezza, fidelity, concurrence e parametri delle disuguaglianze di Bell. Le previsioni analitiche vengono confrontate con i risultati numerici per validare l'implementazione e identificare le diverse firme dell'evoluzione coerente, del dephasing statistico, della depolarizzazione e della PMD.

% Il lavoro teorico e computazionale è completato dallo sviluppo di una piattaforma di laboratorio a due rami per la distribuzione di entanglement di polarizzazione, basata sulla rivelazione di singoli fotoni e su misure di coincidenza, insieme a trasformazioni controllabili della polarizzazione. Il setup rappresenta la controparte sperimentale del framework sviluppato e una base per una futura implementazione di quantum key distribution BBM92 basata sull'entanglement. Nel complesso, la tesi stabilisce un percorso continuo dalla modellizzazione del canale in fibra, attraverso la simulazione numerica, fino alla caratterizzazione sperimentale della distribuzione di entanglement in fibre ottiche.